% 后记（不编号）
\chapter*{\texorpdfstring{\centering 后记}{后记}}
\addcontentsline{toc}{chapter}{后记}
\markboth{后记}{后记}
% 正文后部偶发 \parindent 被冲掉；此处显式恢复中文段首缩进
\setlength{\parindent}{2em}
\setlength{\parskip}{0.2em}

本书的问题核心是\textbf{素数分布}，组织方式是围绕黎曼猜想的\textbf{知识谱系}：
从古典经验公式与素数定理，到 $\zeta$、$\xi$、零点、函数方程与显式公式，再到数值确认与当代前沿简介。
书名将「素数分布」与「黎曼猜想」并列；前言说明：问题核心是素数分布，组织主线是理解该问题与黎曼猜想所需的对象及已证关系。
正文分章依数学依赖关系编排——
先定义与工具，解析延拓、函数方程、零点与显式公式，后为数值计算、几何示意与前沿简介——
以便读者循序核对每一结论所依据的假设与前置结果。

读完全书后，有一点宜特别标明。
第~\ref{ch:riemann_explicit_formula}~章的黎曼显式公式给出 $\psi_0$（及相应计数函数）与 $\zeta$ 非平凡零点之间的\textbf{精确对应}：在表示的意义上，主项与振荡结构均已写出。
然而恒等式本身并不蕴含 $\psi(x)\sim x$。
素数定理所断言的，是在该对应下零点贡献相对于主项为 $o(x)$；这依赖于 $\zeta(1+it)\neq 0$ 等须\textbf{单独建立}的分析事实（第~\ref{ch:prime_number_theorem}~章）。
因此，显式公式没有终结素数定理；素数定理亦非显式公式成立后的直接推论。
宜分清黎曼猜想知识谱系中的不同层面：主项合法性、精确结构、最优误差。

编者对本书所介绍知识谱系可作如下总括：
\textbf{显式公式已从数理上给出「小于 $x$ 的素数个数」的通项型精确结构
（主项与全体非平凡零点贡献之和，已证）；
黎曼猜想只是对截断后（或相对主项）的误差大小给出最优估计
（与 $\operatorname{Re}\rho=1/2$ 等价，未证）。}

本书\textbf{由学习笔记整理而成}：在摘录、推导与图示的基础上，
经整理删改、统一符号与重写叙述后形成现在的章节。
整理时尽量采用可检验的写法，并标明观察与证明的界限；
严格程度与叙述风格难比肩经典文献，疏漏与不完备之处由作者负责，并欢迎指正。
与书中图示、动画相关的 Mathematica、Manim 源程序、视频文件及运行环境说明，
见书后附注与 \url{https://github.com/lichengtong123/lichengtong-RH}。

若本书能帮助读者建立黎曼猜想相关的最小自洽的数学知识体系，
并理清素数分布问题与黎曼猜想之间的关系，
便达到了整理的目的。

\vspace{1.5em}
\begin{flushright}
作者\\
\today
\end{flushright}
